\section{Whole coefficient records on arbitrary carriers}

The derivative ideals along a subset are defined even when the subset
is not globally regular. Their naturality under local biholomorphisms
therefore has an algebraic consequence on such subsets. This extends
the naturality in \dcref{mc:homogenized-along-etale}: regularity and
coherence are unnecessary for the equalities below. All ideals under
consideration have holomorphic members. No assertion of coherence of a
derivative ideal on a singular carrier is made.

\begin{lemma}[Pullback of powers with holomorphic members]
\label{hir:actual-holomorphic-power-pullback}
\uses{mc:homogenized-basic}
\lean{TranscendentalNakamaye.Pending.HironakaActualNonregularCoefficientNaturality.pullback_mul,TranscendentalNakamaye.Pending.HironakaActualNonregularCoefficientNaturality.pullback_idealPow}
\leanok
For a holomorphic map between complex manifolds of arbitrary dimensions,
and an ideal $I$ whose members are holomorphic,
$f^*(I^q)=(f^*I)^q$ for every $q\geq0$.
\end{lemma}
\begin{proof}
The pullback of the unit ideal is the unit ideal: every pullback germ
is holomorphic, and conversely a holomorphic germ is its own coefficient
times the pullback of $1$. This proves the assertion at $q=0$.
Exact product pullback requires only holomorphic members and also holds between different dimensions. Indeed,
composing a finite holomorphic expression in products of source members
gives an expression in products of pullback members. Conversely, write
each of two pullback members as a finite holomorphic combination of
composed source members. Multiplying the expressions and distributing
gives the finite double sum with coefficients the products of the two
coefficients; each generator is the composition of the corresponding
source product. This proves the reverse inclusion. Every member of a
power is holomorphic, since the zeroth power is the holomorphic unit
ideal and every positive power is a holomorphic span of products.
Applying the product identity to $I^q I$ and the induction hypothesis
proves the assertion at $q+1$.
\end{proof}

\begin{lemma}[Pullback of every finite holomorphic sum]
\label{hir:actual-holomorphic-finite-sum-pullback}
\uses{hir:actual-holomorphic-power-pullback}
\lean{TranscendentalNakamaye.Pending.HironakaActualNonregularCoefficientNaturality.pullback_finSumIdeal,TranscendentalNakamaye.Pending.HironakaActualNonregularCoefficientNaturality.pullback_homSum}
\leanok
If each $J_i$ for $i<r$ has holomorphic members, then
$f^*(\sum_{i<r}J_i)=\sum_{i<r}f^*J_i$, including $r=0$. In particular pullback commutes with $\sum_{i<r} A_i B^i$ when the $A_i$ and $B$ have holomorphic members.
\end{lemma}
\begin{proof}
A germ on the left is a finite holomorphic combination of compositions
$a\circ f$, where $a$ is a finite holomorphic combination of members of
the $J_i$ for $i<r$. Composing this latter combination gives a finite
combination of composed individual members, since composition preserves
holomorphicity of its coefficients. Flattening the two finite sums puts
the germ in the right side. Conversely a member of the right side is a
finite combination of members of individual $f^*J_i$. Expand each of
these as a finite combination of composed source members. A source
member of $J_i$ belongs to the source sum by the one-term expression
with coefficient $1$. Thus every expanded generator belongs to the
pullback of the source sum, and flattening proves the required
inclusion. These arguments use no selected index and remain valid for
the empty family: the only expressions then are empty sums and their
holomorphic combinations. Apply this identity to $J_i=A_iB^i$. Product pullback and the preceding power identity identify each summand with $(f^*A_i)(f^*B)^i$, proving the stated particular case.
\end{proof}

\begin{theorem}[Whole homogenized and coefficient records on arbitrary carriers]
\label{hir:actual-nonregular-whole-coefficient-naturality}
\uses{hir:actual-holomorphic-power-pullback,hir:actual-holomorphic-finite-sum-pullback,mc:homogenized-along-etale}
\lean{TranscendentalNakamaye.Pending.HironakaActualNonregularCoefficientNaturality.pullback_homMarkedAlong,TranscendentalNakamaye.Pending.HironakaActualNonregularCoefficientNaturality.pullback_coeffMarkedAlong,TranscendentalNakamaye.Pending.HironakaActualNonregularCoefficientNaturality.pullback_coeff_homMarkedAlong}
\leanok
Let $M=(N,E,I,d)$ have arbitrary carrier $N$, holomorphic ideal members,
and any natural mark $d$. For a holomorphic local biholomorphism $f$,
the whole homogenized record and the whole coefficient record commute
with pullback. Neither global regularity of $N$, coherence, nor positive
mark is required.
\end{theorem}
\begin{proof}
The derivative identity along an arbitrary subset is
$f^*D_N^iI=D_{f^{-1}N}^i(f^*I)$. Every member of each such derivative
ideal is holomorphic. Put $T=D_N^{d-1}I$. Homogenization is the finite
sum of $D_N^iI\,T^i$ for $i<d$. Product pullback, power pullback, and
finite-sum pullback identify this whole ideal with the corresponding
homogenized ideal on $f^{-1}N$. The coefficient ideal is the finite sum
of the prescribed marked-sum powers of $D_N^iI$ for $i<d$. The same
power and finite-sum identities identify its whole ideal, with the
identical exponents because the mark is unchanged by pullback. For
$d=0$ both finite sums are empty and the preceding finite-sum identity
still applies. In both constructions the carrier is $f^{-1}N$, the
boundary is the pullback of the same entire positioned list, and the
output mark is computed from the same original marks. Equality of
these fields and of the whole ideals proves equality of the whole
marked records. Every member of the homogenized ideal is holomorphic because it is a finite holomorphic span. Applying the coefficient identity to this homogenized record and then the homogenization identity proves the whole composite identity for $\mathcal C_N(\mathcal H_N(M))$.
\end{proof}
